\section{Other tested variants}\label{sec:variants}

\subsection{The three-shift digamma construction}

Before the mixed moments, \rec{R15} adapted Litsey's divided-difference
construction and rational-antidifference argument
\cite[Sections~3--6]{Litsey}. It gives an all-slope nonvanishing result,
although the arithmetic and real estimates did not close in the tests.
Let $\psi=\Gamma'/\Gamma$ be the digamma function. For rational
$\lambda$, put
\[
 F(z)=\psi(z+1/4)+(24\lambda-5/2)\psi(z+1/2)
                       +(3/2-24\lambda)\psi(z+1).
\]
The weights sum to zero, and $F'(0)=8(G+\lambda\pi^2)$. Its increment
\[
 \Delta(z)=F(z+1)-F(z)=\frac1{z+1/4}
 +\frac{24\lambda-5/2}{z+1/2}+
 \frac{3/2-24\lambda}{z+1}
\]
is rational. At nonnegative integer nodes define
$f_t=\sum_{j=0}^{t-1}\Delta(j)$ and
$d_t=\sum_{j=0}^{t-1}\Delta'(j)$. Choose a row-node set $R$ and a
column-node set $C$, each of size $N=r+m$, whose union is
$\{0,\ldots,r+2m-1\}$ and whose intersection is a consecutive block
of $r\ge1$ nodes. The matrix entries, for $t\in R$ and $u\in C$, are
\[
 M(t,u;X)=\frac{f_u-f_t}{u-t}\quad(t\ne u),\qquad
 M(t,t;X)=X+d_t.
\]
At $X=F'(0)=8(G+\lambda\pi^2)$, these entries are the divided
differences of $F$, with $F'(t)$ on a common-node diagonal.
The determinant is a rational polynomial in this single target.
The following proposition supplies nonvanishing for the determinant
together with its common-node cofactors.

\begin{proposition}\label{prop:digamma-rank}
For every complex $X$, this square matrix has corank at most one.
If it is singular, some common-node diagonal cofactor is nonzero.
Thus the determinant and these cofactors have no common complex zero.
\end{proposition}
\begin{proof}
Let $Q_C(z)=\prod_{u\in C}(z-u)$. A right kernel vector corresponds to
polynomials $U,W$ of degree below $N$ through
\[
 U=Q_C\sum_{u\in C}\frac{v_u}{z-u},\qquad
 W=Q_C\sum_{u\in C}\frac{v_uf_u}{z-u}.
\]
The kernel equations are precisely $W(t)=f_tU(t)$ at every union
node and $W'(t)=f_tU'(t)+(X+d_t)U(t)$ at each common node.
The transpose has the same interpolation conditions.

For two such pairs, $W_1U_2-W_2U_1$ has $2N$ zeros counted with
multiplicity and degree at most $2N-2$, so it vanishes. All pairs
therefore have the same reduced ratio. If the kernel dimension is
$d$, canceling the highest $d-1$ coefficients of their polynomial
multipliers gives a pair of maximum degree at most $N-d$.
Write $\Delta=n_\Delta/d_\Delta$ in lowest terms, with denominator
degree $q\le3$. The zero sum of the weights gives
$\deg n_\Delta\le q-2$. For a pair of maximum degree $D$, form
\[
 T(z)=d_\Delta(z)\{W(z+1)U(z)-W(z)U(z+1)\}
                         -n_\Delta(z)U(z)U(z+1).
\]
Its degree is at most $2D+q-2$. It vanishes at every adjacent pair
of union nodes and has a double zero at each adjacent pair of common
nodes, giving $2N-2$ zeros. If $d\ge2$, its degree is at most
$2N-3$, so $T=0$. This would make $W/U$ a rational antidifference
of $\Delta$. But a rational difference has residue sum zero in each
pole class modulo $\Z$, whereas the class of $-1/4$ contains one
pole of $\Delta$ with residue one. Hence $d\le1$.

If $U$ vanished at all common nodes, so would $W$. Dividing by their
common-node product leaves degrees at most $m-1$ and preserves the
value conditions at all union nodes. If $m=0$, the reduced pair is
already zero. For $m\ge1$, the analogous polynomial $T$ has degree at most $2m-1$ but at least $r+2m-1\ge2m$ zeros,
again impossible. Some common node therefore has $U(t)\ne0$.
Both the left and right kernel coordinates there are nonzero,
so the corresponding diagonal adjugate entry is nonzero.
\end{proof}

For $\lambda=1$ and border ratio $2m/r=1/5$, put
$P(X)=\det M(X)$ and $\xi=8(G+\pi^2)$. Let
$K=1/\cont(P)>0$, so $KP$ is the complete primitive integer polynomial.
The following values are rounded from certified finite enclosures;
every logarithm is divided by $r^2$:
\begin{center}
\begin{tabular}{rrrr}
\toprule
$r$&Multiplier cost&Raw decay&Primitive growth\\
&$\log K/r^2$&$-\log|P(\xi)|/r^2$&$\log|KP(\xi)|/r^2$\\
\midrule
20&3.384060&3.168234&.215826\\
40&3.675448&3.417346&.258102\\
80&3.771926&3.541927&.229999\\
160&3.813989&3.604388&.209601\\
\bottomrule
\end{tabular}
\end{center}
The last column is the first minus the second. Its positive entries
mean that the primitive values have absolute value greater than one at
these sizes, after including the full coefficient gcd.
The slope $\lambda=1/16$ is a warning against
extrapolation: its primitive rates are negative at $r=20,40,80$,
but positive at $r=160$. Litsey's special slopes serve as controls;
their arithmetic cancellation cannot be assumed at every rational slope.

\subsection{Two further Pad\'e variants}

The reassessment \rec{R28} tested two changes to the family of
Section~\ref{sec:pade}. First, writing $X=t(t-1)$ converts powers
$X^j$ into the even factorial basis through
\[
 X^j=\sum_\ell S(j,\ell)T_{2\ell}(t),\qquad
 S(j+1,\ell)=S(j,\ell-1)+\ell(\ell+1)S(j,\ell).
\]
The first nontrivial Hankel determinant of the resulting remainders is
$E_0(E_2+2E_1)-E_1^2$. Nine tests constructed the entire polynomial
in $\beta(4)$ and reduced it to primitive integer coefficients.
At seed index $n=12$, determinant sizes $1,2,4,6$ gave logarithms
of absolute primitive values approximately $78.77,164.19,324.32,508.39$.
The coefficient calculations are exact; these decimal evaluations
were exploratory. No estimate uniform in growing determinant size
was obtained.

Second, lowering the numerator multiplicity and retaining the crossed
rational residues gives the real integral
\[
 J_{n,m}=\frac2{(n!)^{m-2}}\int_0^\infty
 \frac{y\mu'''(y)\prod_{j=1}^{n/2}(y^2+(j-1/2)^2)^m}
      {\prod_{j=0}^n(y^2+j^2)}\dd y,
 \qquad \mu(y)=\frac1{2\cosh\pi y}.
\]
Partial fractions express it in $\Q+\Q\beta(4)$. The $m=5$ case
agrees exactly with the earlier construction at the four comparison
indices. For $m=3,n=2$, its primitive form is
$90\beta(4)-89$, of absolute value about $.00501$. At $n=4,8,12,20$,
the corresponding logarithms are approximately $12.50,37.93,61.43,121.56$.
Lower multiplicity reduces a cost, but these tests supply no favorable
asymptotic comparison.
